% TOP_PROBLEM: {"id": 30007432, "title": "Nearly minimal Littlewood maxima with positive minima", "category_id": 8, "category": {"id": 8, "name": "analysis", "display_name": "Analysis", "description": "Limits, continuity, calculus, and function theory.", "slug": "analysis", "order_index": 8, "created_at": "2026-07-31T15:26:25.671Z"}, "status": "open"}
% FIRST_POSTED: 2026-10-06
% ATTEMPT_STATUS: solved
% OPENAI_CLAIM: {"schema_version":1,"source_repo":"https://github.com/openai/math","source_commit":"adc7f1241b42e322a6451854ab7e4b4c146bf78a","release_date":"2026-10-06","match":"full","resolution":"proved","independently_reviewed":false,"second_pass":"agreed","families":[{"family":"076","title":"Real ultraflat Littlewood polynomials and unbounded binary merit factors","subject":"Real and complex analysis","manuscripts":[{"title":"Nearly minimal maxima and positive minima of Littlewood polynomials","pdf_path":"preprints/Nearly-minimal-maxima-and-positive-minima-of-Littlewood-polynomials-October-5-2026/littlewood-lower-envelope.pdf","readme_path":"preprints/Nearly-minimal-maxima-and-positive-minima-of-Littlewood-polynomials-October-5-2026/README.md","date":"2026-10-05","theorem_ref":"thm:main (sections/introduction.tex:18)","statement_sources":[{"path":"preprints/Nearly-minimal-maxima-and-positive-minima-of-Littlewood-polynomials-October-5-2026/build/sections/introduction.tex","line":18,"label":"thm:main"}]}],"lean":{"doc_path":"lean/docs/076.md","comparators":[],"declarations":[],"covers_main_theorem":false},"reasoning_trace":null}]}
\documentclass[11pt]{article}
\usepackage[margin=25mm]{geometry}
\usepackage{fontspec}
\setmainfont{Latin Modern Roman}
\usepackage{amsmath,amssymb,mathtools}
\usepackage{unicode-math}
\setmathfont{Latin Modern Math}
\usepackage{enumitem}
\usepackage{xurl,hyperref,bookmark}
\hypersetup{colorlinks=true,urlcolor=blue}
\setcounter{secnumdepth}{0}
\setlength{\parindent}{0pt}
\setlength{\parskip}{5pt}
\setlength{\emergencystretch}{3em}
\begin{document}
\section*{Nearly minimal Littlewood maxima with positive minima}

\subsection{OpenAI's claim}
\textit{Littlewood polynomials with a lower margin}: The paper claims sign polynomials at every sufficiently large length with maximum at most $(1+\eta)\sqrt N$ and minimum at least $\sqrt N/16$, for any fixed positive $\eta$. Both inequalities concern the same polynomial.

\subsection{Mathematical statement of the claimed result}
\paragraph{Mathematical claim 1.} Given $\eta>0$, there is $N_0(\eta)$ such that for every integer $N\ge N_0(\eta)$ one can choose real signs $\epsilon_0,\ldots,\epsilon_{N-1}$ with \[\frac{\sqrt N}{16}\le\left|\sum_{k=0}^{N-1}\epsilon_kz^k\right|\le(1+\eta)\sqrt N\quad(|z|=1).\] The same polynomial has both bounds uniformly on the circle; the lower constant is $1/16$.

\subsection{Outline of the argument}
\textit{Littlewood polynomials with a lower margin}: The relaxed-coefficient construction for small maxima is augmented on arcs where its leading amplitude is too small. Carefully matched quadratic phases create a positive lower margin without introducing excessive Fourier coefficients outside the permitted degree range. A rounding estimate expressed in terms of the remaining coefficient defect preserves the margin while turning the coefficients into signs.

\subsection{Scope relative to this problem}
\textit{Littlewood polynomials with a lower margin}: The primary statement says: For every eta>0 and every sufficiently large integer length N, a real sign polynomial has modulus between sqrt(N)/16 and (1+eta)sqrt(N) everywhere on the circle. The target retains these quantified objects and hypotheses, so its positive assertion is covered.

With the hypotheses and conventions above, the claimed result addresses the whole mathematical assertion on this page.

\subsection{Formal verification}
Family 076 describes Lean coverage for near-minimal upper maxima and convergence of finite-exponent modulus means; the later uniform lower-envelope conclusions are separate.

\textit{Littlewood polynomials with a lower margin}: No manuscript-specific Comparator is listed in the pinned inventory.

The documented selection does not cover the entire manuscript assertion relevant here.

This site has not run the Lean development. These links describe OpenAI’s selected formal statements and should be read with their stated scope.

\subsection{Status caveat}
OpenAI reports that an internal model produced the manuscripts. Their arguments have not been independently verified as proofs by this site. Any portion outside the documented formal scope may contain mistakes, and the catalogue label reports the announced mathematical claim.

\subsection{Sources}
\begin{itemize}
\item OpenAI, \emph{Nearly minimal maxima and positive minima of Littlewood polynomials}, 2026-10-05; citation key \texttt{OAI:Nearly-minimal-maxima-and-positive-minima-of-Littlewood-polynomials-October-5-2026}. \href{https://github.com/openai/math/blob/adc7f1241b42e322a6451854ab7e4b4c146bf78a/preprints/Nearly-minimal-maxima-and-positive-minima-of-Littlewood-polynomials-October-5-2026/littlewood-lower-envelope.pdf}{PDF (GitHub)}; \href{https://github.com/openai/math/raw/adc7f1241b42e322a6451854ab7e4b4c146bf78a/preprints/Nearly-minimal-maxima-and-positive-minima-of-Littlewood-polynomials-October-5-2026/littlewood-lower-envelope.pdf}{PDF (direct)}; \href{https://github.com/openai/math/blob/adc7f1241b42e322a6451854ab7e4b4c146bf78a/preprints/Nearly-minimal-maxima-and-positive-minima-of-Littlewood-polynomials-October-5-2026/README.md}{Manuscript page and citation}; \href{https://github.com/openai/math/blob/main/preprints/Nearly-minimal-maxima-and-positive-minima-of-Littlewood-polynomials-October-5-2026/littlewood-lower-envelope.pdf}{Latest version}; Statement source: \href{https://github.com/openai/math/blob/adc7f1241b42e322a6451854ab7e4b4c146bf78a/preprints/Nearly-minimal-maxima-and-positive-minima-of-Littlewood-polynomials-October-5-2026/build/sections/introduction.tex\#L18}{thm:main}.
\item Family 076: \href{https://github.com/openai/math/blob/adc7f1241b42e322a6451854ab7e4b4c146bf78a/lean/docs/076.md}{Lean scope}.
\item \href{https://github.com/openai/math/blob/adc7f1241b42e322a6451854ab7e4b4c146bf78a/CONTENTS.md}{OpenAI catalogue entry}.
\end{itemize}
\end{document}
